\section{\ref{sec:dofa}장. \nt{DOPA}를 이용한 학습}

시냅스의 메커니즘(꾸러미, 동시성 검출기, 칼슘, 가소성 창)과 예측 오차로서의 도파민의 행동은 직접 측정되었다. 규칙들이 기울기로 이어진다는 것과 브로드캐스트의 대가는 정리다. 선택 학습의 결과는 이 책의 모델로 계산한 것이다. 오차를 계산하는 회로는 가설이다.

{\footnotesize
\setlength{\tabcolsep}{3pt}\renewcommand{\arraystretch}{1.15}
\begin{longtable}{>{\raggedright}p{0.5cm}>{\raggedright}p{4.0cm}>{\raggedright}p{5.6cm}>{\raggedright}p{2.3cm}>{\raggedright\arraybackslash}p{1.7cm}}
\toprule
No. & 명제 & 실험과 측정한 것 & 출처 & 상태 \\
\midrule\endfirsthead
\toprule
No. & 명제 & 실험과 측정한 것 & 출처 & 상태 \\
\midrule\endhead
1 & 인공 신경망과 같은 형태의 오차 역전파는 뇌에서 발견되지 않았다 & 직접적인 실험은 없다. 이것은 부재에 관한 명제다. 리뷰들은 이 알고리즘에 무엇이 필요한지(순방향 연결을 되풀이하는 역방향 연결, 별도의 단계)와 어떤 생물학적 근사가 제안되었는지를 다룬다. & \cite{Lillicrap2020, WhittingtonBogacz2019} & 가설 \\
2 & 가중치는 방출 부위 수, 방출 확률, 한 꾸러미에 대한 응답으로 분해된다: $w=N_s\,p\,A_1$~\eqref{eq:quantal} & 방출을 낮춘 개구리 신경근 시냅스: 수신자의 응답은 꾸러미 하나에 대한 자발적 미니 응답의 정수배로 계단식으로 변한다. 스파이크당 꾸러미 수는 푸아송 법칙을 따른다. & \cite{delCastillo1954} & 측정됨 \\
3 & 수신자의 수용체는 동시성 검출기이고, 칼슘이 가중치 변화를 일으킨다 & 생쥐 뉴런의 패치 클램프: 휴지 전위에서는 Mg$^{2+}$ 이온이 글루탐산으로 열린 채널을 막고, 탈분극되면 빠져나간다. 해마 절편: 수신자 안의 칼슘 킬레이터는 장기 강화를 막고, 세포 안에서 빛으로 칼슘을 풀어 주면 그것만으로 전달이 강화된다. & \cite{Nowak1984, Malenka1988calcium} & 측정됨 \\
4 & 결정은 어느 쪽이든 실행한다: 수신자에서는 $A_1$이 커지고, 송신자에서는 $p$가 바뀐다 & 가시 하나 위에서 빛으로 풀어 준 글루탐산은 가시와 그 수용체를 흐르는 전류를 빠르고 지속적으로 키운다. 강화에는 AMPA 수용체의 삽입이 따른다. 수신자의 탈분극은 송신자의 방출을 한동안 억누른다. 이 효과는 틈을 거슬러 가는 엔도카나비노이드가 전달한다(\nt{GABA} 입력에서 확인됨). 단기 고갈은 송신자의 방출 확률에 달려 있다. & \cite{Matsuzaki2004, Malinow2002, WilsonNicoll2001, Tsodyks1997} & 측정됨 \\
5 & 송신자와 수신자가 함께 활동하면 시냅스가 강해진다~\eqref{eq:hebb} & 마취한 토끼: 해마 입력 경로에 짧은 스파이크 열을 준 뒤 응답이 $30$~min에서 $10$~h까지 강화된다. 해마 절편에서 수신자의 스파이크는 같은 순간 활동한 시냅스만 강화한다. & \cite{Bliss1973LTP, Kelso1986hebb} & 측정됨 \\
6 & 규칙~\eqref{eq:oja}은 입력 상관행렬의 주 고유벡터를 찾는다(명제~\ref{prop:oja}) & 정리이며 증명은 본문에 있다. 뉴런이 자기 입력의 주성분을 배운 실험은 없다. 시냅스에서 가중치 증가를 제한하는 메커니즘은 밝혀지지 않았다. & \cite{Oja1982} & 이론 \\
7 & 시각에 따른 헤브 규칙: 창은 약 $\pm20\ms$(그림~\ref{fig:stdp}). 반복되는 순서열에서 사슬이 자란다 & 배양한 해마 뉴런 쌍: 송신자의 스파이크 뒤 $20\ms$ 안의 수신자 스파이크는 지속적 강화를, $20\ms$ 앞의 스파이크는 약화를 준다. 둘 다 NMDA 수용체에 의존한다. 그림의 비율 $A_-=0.6A_+$는 예시다. 이렇게 신경망에서 사슬이 자란다는 것은 직접 측정되지 않았다. & \cite{BiPoo1998} & 측정됨 \\
8 & 적격 흔적~\eqref{eq:threefactor}: 함께 일어난 활동 뒤 $1$--$2\,\text{s}$ 안에 온 도파민은 연결을 강화하고, 더 늦으면 그렇지 않다(명제~\ref{prop:threefactor}) & 선조체 절편, 글루탐산 입력과 도파민 입력의 개별 광유전학 자극: 도파민이 글루탐산 입력 뒤 $0.3$--$2\,\text{s}$ 안에 올 때만 가시가 자란다. 창은 cAMP의 빠른 상승과 분해로 정해진다. 피질에서는 짝짓기가 강화와 약화에 대해 따로따로 흔적을 남기고, 모노아민 수용체가 이것을 몇 초 규모의 창 안에서 장기 변화로 바꾼다. & \cite{Yagishita2014, He2015eligibility} & 측정됨 \\
9 & 몇 초 길이의 창은 도파민 없이도 있다: 세 번째 신호는 수신자의 강한 흥분이다 & 살아 있는 생쥐, CA1 영역: 수신자의 플래토 사건 하나가 그 앞뒤 몇 초 안에 온 입력을 강화하고, 한 번 지나가는 것만으로 장소장을 만든다. & \cite{Bittner2017} & 측정됨 \\
10 & 3요소 규칙의 평균 걸음은 보상의 기울기를 따른다~\eqref{eq:perturbation}(명제~\ref{prop:gradient}) & 무작위 이탈로 학습할 때의 기울기에 관한 정리. 본문에서는 작은 잡음에 대해 유도했다. & \cite{Williams1992} & 이론 \\
11 & 잡음은 탐색이다: 신경망은 자기의 무작위 이탈로부터 배운다 & 어린 금화조: 기저핵 회로의 출력 핵을 끄면 노래의 변동성이 급격하고 가역적으로 사라진다. 성체 금화조: 음높이가 평균의 한쪽에 있는 음절을 부를 때 방해음을 주면 새는 음높이를 빠르게 반대쪽으로 옮긴다. 자기 산포로부터 배우는 것이다. & \cite{Olveczky2005, TumerBrainard2007} & 측정됨 \\
12 & 두 선택지 중 고르기는 $\tau_e=1\,\text{s}$, $\delta=R-\bar R$에서 학습되고, $\tau_e=50\ms$에서는 학습되지 않는다. 기대값을 빼지 않으면 가중치가 한없이 자라고, 학습률이 크면 신경망이 더 나쁜 선택을 굳힐 수 있다 & \texttt{models/dofa\_learning.py}, $\eta=0.05$, 시도 $500$회, 시드 $6$개. $\tau_e=1\,\text{s}$: $A$ 선택 비율 $0.65$--$0.79\to0.96$--$1.00$, 가중치 $0.85$--$1.33$. $\tau_e=50\ms$: $0.38$--$0.55$, 가중치 변화 없음. $\delta=R$: 승자의 가중치 $860$--$1190$. $\delta=R$, $\eta=0.5$, 시드 $3$개: 하나가 $B$에 고정($0.04\to0$). 스크립트가 네 경우를 모두 출력하며, 재계산 결과는 본문과 일치했다. & \texttt{models/\allowbreak dofa\_\allowbreak learning.py} & 모델 \\
13 & 하나의 공통 신호로 학습하는 시간은 잡음을 내는 뉴런 수에 비례해 늘어난다(명제~\ref{prop:broadcastcost}) & 선형 신경망의 학습 곡선: 뉴런 섭동은 정확한 기울기보다 출력 뉴런 수만큼 느리고, 가중치 섭동은 거기에 입력 수만큼 더 느리다. & \cite{Werfel2005} & 이론 \\
14 & \nt{DOPA}의 출력은 무엇보다 행동 선택 영역으로 간다 & 쥐의 흑질선조체 \nt{DOPA} 뉴런 개개의 축삭 완전 재구성: 세포 하나의 가지가 선조체 부피의 $0.45$--$5.7\,\%$(평균 $2.7\,\%$)를 덮고, 선조체 밖에는 가지가 거의 없다. & \cite{Matsuda2009} & 측정됨 \\
15 & 피질은 자기 입력을 예측하도록 배우며, 오차는 그 자리에서 계산된다 & 리뷰: 감각 피질에는 기대한 입력과 실제 입력의 불일치에 대한 응답이 있다. 이것이 피질 학습의 일반 규칙이라는 것은 증명되지 않았다. & \cite{Keller2018} & 가설 \\
16 & 도파민은 오차~\eqref{eq:ctd}처럼 행동한다: 급증, 예고 신호로의 이동, 급감(그림~\ref{fig:ctdcases}) & 원숭이의 \nt{DOPA} 뉴런: 예상 못 한 주스에 급증. 학습 뒤에는 예고 신호에 급증이 생기고 약속된 주스에는 응답이 없다. 주스가 오지 않으면 급감. 연속 형태~\eqref{eq:ctd}는 이론이다. & \cite{Schultz1997, Doya2000} & 측정됨 \\
17 & $\dd V/\dd t$를 계산하고 기대를 빼는 회로 & 생쥐, 배쪽 피개 영역의 기록과 광유전학: \nt{DOPA} 뉴런은 보상에서 기대를 뺀다. 이웃한 \nt{GABA} 뉴런은 보상이 기대될 때 이들을 억제하며, 이 뉴런들을 흥분시키거나 억제하면 오차가 바뀐다. 도함수가 어떻게 계산되는지는 밝혀지지 않았다. & \cite{Eshel2015arithmetic} & 가설 \\
18 & \nt{DOPA} 뉴런은 페이스메이커이며, 급감은 급증보다 얕다 & 절편에서 \nt{DOPA} 뉴런은 느린 페이스메이커 전위 위에서 자발적으로, 매우 규칙적으로 발화한다. 원숭이에서 발화 빈도는 양의 오차를 정량적으로 따르지만 음의 오차는 약하게만 전달한다. & \cite{GraceOnn1989, Bayer2005} & 측정됨 \\
19 & 행위자와 비평자(그림~\ref{fig:actorcritic}) & 알고리즘은 강화 학습 이론에서 왔다. 사람(fMRI)에서 배쪽 선조체는 선택이 있든 없든 비평자처럼 행동하고, 등쪽 선조체는 행동을 골라야 할 때만 그렇다. & \cite{SuttonBarto2018, ODoherty2004actor} & 가설 \\
20 & 두 번째 계열의 수용체를 가진 뉴런에서는 규칙의 $\delta$ 부호가 뒤집힌다 & 선조체 절편: D1 수용체를 가진 뉴런에서는 짝짓기가 D1이 필요한 강화를, D2를 가진 뉴런에서는 D2가 필요한 약화를 준다. & \cite{Shen2008} & 측정됨 \\
\bottomrule
\end{longtable}}
